% ============================================================
%  PART V --- Settings
% ============================================================
\part{Settings}

% ── Chapter 20: Settings Overview ───────────────────────────
\chapter{Paramètres --- Aperçu général}
\label{ch:settings}

Open Settings by tapping \iconlabel{settings-iec}{Settings} in the Bottom Bar.
La configuration de bureau et iOS est organisée sur dix onglets. Android ajoute un
eleventh \textbf{Android} tab after System.
Les modifications prennent effet immédiatement à moins qu'un redémarrage ne soit explicitement noté.

\screenshot{settings-overview}{Panneau de configuration avec les onglets appropriés à la plate-forme visibles dans la barre des onglets}{}

\rowcolors{2}{LtGray}{white}
\begin{longtable}{C{0.6cm}L{3.0cm}L{6.8cm}L{3.3cm}}
\toprule
\rowcolor{Navy}
\color{white}\sffamily\bfseries \# &
\color{white}\sffamily\bfseries Tab name &
\color{white}\sffamily\bfseries Primary function &
\color{white}\sffamily\bfseries Restart needed? \\
\midrule
0 & Main           & Window mode, language, UI scale, comfort preset, launcher grid, coordinate format & Language and virtual keyboard only \\
1 & Top Bar        & Drag-and-drop layout of the Top Bar instrument strip & No \\
2 & Bottom Bar     & Drag-and-drop layout of the Bottom Bar action strip & No \\
3 & Widgets        & Define and edit reusable data widget definitions & No \\
4 & Comfort        & Select and customise visual comfort presets & No \\
5 & Connection     & Signal~K server URL, authentication token, mDNS discovery & No \\
6 & Signal~K       & Map semantic names to actual Signal~K paths & No \\
7 & Units          & Choose display units per measurement category & No \\
8 & Applications   & Manage the launcher: pages, tiles, folders, app details & No \\
9 & System         & Diagnostics, logging, performance, configuration management & No (some actions restart) \\
10 & Android       & Optional Home role, native launcher application selection, soft navigation controls & No (Android only) \\
\bottomrule
\end{longtable}

On Android, the Android tab lists only launchable \code{MAIN + LAUNCHER}
activités. Sélectionner une application l'ajoute à la palette des Applications partagées sans
demander une large visibilité du paquet. Identité de l'application autochtone, icône, paquet et
l'activité sont gérées par Android. Désélection d'une application supprime sa page de lancement
Références. FairWindSK ne devient jamais l'application Home par défaut automatiquement, et la
l'intention du lanceur ordinaire reste disponible. Contrôles Soft Back, Home et Recents
apparaissent seulement lorsque Android signale que les clés matérielles correspondantes sont absentes;
Les récents ne contiennent que des activités autochtones lancées par FairWindSK.

% ── Chapter 21: Tab 0 Main ──────────────────────────────────
\chapter{Paramètres --- Principale}
\label{ch:settings-main}
\settingstab{0}{Main}

L'onglet principal contrôle le comportement fondamental de l'exécution : mode d'affichage de la fenêtre, langue, élément UI
taille, confort actif préréglé, dimensions de grille de lancement et format de coordonnées.

\screenshot{settings-main}{Paramètres --- Onglet principal: mode fenêtre, langue, échelle d'interface utilisateur, préréglage confort, grille de lancement, format de coordonnées}{}

\section{Mode fenêtre}

\rowcolors{2}{LtGray}{white}
\begin{longtable}{L{2.4cm}L{5.6cm}L{5.7cm}}
\toprule
\rowcolor{Navy}
\color{white}\sffamily\bfseries Mode &
\color{white}\sffamily\bfseries Behaviour &
\color{white}\sffamily\bfseries Notes \\
\midrule
Fenêtre&Fenêtre de bureau redimensionnable à gauche/haut/largeur/hauteur configurée.&
  Development, testing, multi-monitor setups. \\
Centred&Fenêtre centrée sur la largeur et la hauteur configurées. Champs de position désactivés.&
  Preferred windowed mode without an exact position. \\
Maximisé&Espace de travail complet.&
Sur Raspberry Pi~OS, utilise la géométrie de la zone de travail de l'écran directement,
  bypassing the window manager. \\
Plein écran&Fenêtre de kiosque sans cadre, toujours en haut.&
Recommandé pour les présentoirs de barre dédiés.
  On labwc desktops, set \code{autohide=true}
  in \code{wf-panel-pi.ini} if the panel overlays FairWindSK. \\
\bottomrule
\end{longtable}

\section{Langue}

FairWindSK est livré avec quatre langages d'interface entièrement pris en charge.
Le réglage de la langue affecte tous les textes natifs du shell --- étiquettes de barre, pages de paramètres,
drawers, warnings, and status messages --- as well as the \code{Accept-Language}
header envoyé au contenu web intégré.

\begin{itemize}
  \item \textbf{System default} --- follows the OS locale when FairWindSK supports
cette langue; revient à l'anglais pour chaque langue OS non soutenue.
  \item \textbf{English} --- forces English regardless of OS locale.
  \item \textbf{Français} --- forces French.
  \item \textbf{Español} --- forces Spanish.
  \item \textbf{Italiano} --- forces Italian. All native shell text is displayed in Italian.
\end{itemize}

\subsection{Comment fonctionne la localisation}

The \code{Localization} module normalises the selected language, chooses the
effective \code{QLocale}, and builds the web \code{Accept-Language} header.
The Qt translator (\code{.qm} resource) is loaded during startup, so native widgets
et les vues Web intégrées commencent avec la même langue et la même culture.
Translation source files live in \code{resources/i18n/} (for example,
\code{fairwindsk\_fr.ts}, \code{fairwindsk\_es.ts}, and
\code{fairwindsk\_it.ts}).

\subsection{Ajouter ou améliorer une langue}

See \code{docs/multilingual.md} in the repository for the full contributor workflow:
adding translatable strings with \code{tr(\ldots)}, updating \code{.ts} files with
Qt Linguist, registering a new language in \code{Localization.cpp} and
\code{ui/settings/Main.cpp}, and updating \code{CMakeLists.txt}.

\begin{warningbox}
Un redémarrage est nécessaire après avoir changé la langue.
FairWindSK montre un tiroir d'information lorsque la langue est changée.
Sans redémarrage, les widgets Qt natifs et la vue Web intégrée peuvent utiliser différents
langues.
\end{warningbox}

\section{Préréglage de l'échelle d'assurance-chômage}

Quatre préréglages : Petit, Normal (par défaut), Grand, Extra Large.
When \key{Auto Scale} is ticked, FairWindSK selects the preset based on display resolution
et DPI.

\section{Préréglage de la vue de confort et commutation automatique}

Sélectionne le confort actif directement dans l'onglet principal.
When \key{Auto Comfort View} is enabled, FairWindSK switches presets automatically
using Signal~K sun-position data (\skpath{environment.sun}).
Les deux conditions doivent être remplies pour que le mode automatique soit disponible:

\begin{enumerate}
  \item The \skpath{environment.sun} path must be configured in the Signal~K tab.
  \item The Signal~K server must be connected and delivering values on that path.
\end{enumerate}

\section{Grille de lancement}

Deux boîtes de spin définissent le nombre de lignes et de colonnes de tuiles app par page de lanceur.
Moins de rangées et de colonnes produisent des tuiles plus grandes qui sont plus faciles à taper.
Plus produire plus d'applications visibles par page.

\section{Format des coordonnées}

Trois formats sont disponibles dans toute l'interface (Widget de position de la barre supérieure, barre POB,
Barre d'ancrage, Mes points de coordonnées :

\rowcolors{2}{LtGray}{white}
\begin{longtable}{L{3.2cm}L{4.8cm}L{5.7cm}}
\toprule
\rowcolor{Navy}
\color{white}\sffamily\bfseries Format &
\color{white}\sffamily\bfseries Example &
\color{white}\sffamily\bfseries When to use \\
\midrule
DD.DDDDD\textdegree{} & 48.85341\textdegree{} N & GPS logs, digital charting software. \\
DD\textdegree{}MM.MMM' & 48\textdegree{}51.205' N & Standard nautical format. Recommended for most sailors. \\
DD\textdegree{}MM'SS.S'' & 48\textdegree{}51'12.3'' N & Traditional celestial navigation format. \\
\bottomrule
\end{longtable}

\section{Clavier virtuel}

Active le clavier virtuel Qt pour les déploiements à écran tactile seulement.
Un redémarrage est nécessaire après le basculement.
Sur Raspberry~Pi~OS, le script d'aide au démarrage applique la variable d'environnement correcte
(\code{QT\_IM\_MODULE=qtvirtualkeyboard}) automatically.

% ── Chapter 22: Tabs 1 and 2 --- Bar Editors ──────────────────
\chapter{Paramètres --- Éditeurs de barre supérieure et de barre inférieure}
\label{ch:settings-bars}
\settingstab{1}{Top Bar} \settingstab{2}{Bottom Bar}

Les onglets ~1 et ~2 partagent le même éditeur visuel de glisser-déposer --- un pour le Top~Bar et un
pour le Bas-Bar.

\screenshot{settings-topbar-editor}{Paramètres --- Éditeur de barre supérieure : bande d'aperçu en direct, barre d'outils d'édition et palette de widget}{}

\section{Bande de prévisualisation en direct}

L'aperçu en haut rend une simulation complète de la barre en utilisant le confort actuel
couleurs prédéfinies et données réelles de widget.
Appuyez sur n'importe quel élément de l'aperçu pour le sélectionner.

\section{Modifier les boutons de barre d'outils}

These are the exact icons defined in \code{BarLayoutSettings.cpp}.

\rowcolors{2}{LtGray}{white}
\begin{longtable}{C{1.5cm}L{4.0cm}L{8.2cm}}
\toprule
\rowcolor{Navy}
\color{white}\sffamily\bfseries Icon &
\color{white}\sffamily\bfseries Button &
\color{white}\sffamily\bfseries Action \\
\midrule
\iconlg{layout-horizontal-minimize} & Minimize width  &
  Widget uses only the space its content requires. \\
\iconlg{layout-horizontal-maximize} & Maximize width  &
  Widget expands to fill all available space horizontally. \\
\iconlg{layout-vertical-minimize}   & Minimize height &
  Widget uses the standard bar height. \\
\iconlg{layout-vertical-maximize}   & Maximize height &
  Widget uses the full bar height (both header and value rows fully visible). \\
\iconlg{arrow-left-google}          & Move left       &
  Moves the selected widget one position to the left. \\
\iconlg{arrow-right-google}         & Move right      &
  Moves the selected widget one position to the right. \\
\iconlg{delete-google}              & Remove selected &
  Returns the selected widget to the palette. \\
\iconlg{refresh-google}             & Reset defaults  &
  Restores the bar to its built-in default layout. \\
\bottomrule
\end{longtable}

\section{Afficher les toggles}

Avec un widget sélectionné, quatre cases contrôlent ce qui est montré :
\key{Show Icon}, \key{Show Text} (label), \key{Show Units}, \key{Show Trend} (arrow).

\section{Palette Widget}

La palette ci-dessous liste tous les éléments disponibles.
\textbf{Data widgets} are backed by Signal~K paths.
\textbf{Fixed shell controls} --- Apps, POB, Autopilot, Anchor, Alarms, My~Data, Settings,
L'état Signal~K, Horloge --- peut être déplacé entre la barre supérieure et la barre inférieure.
\textbf{Separators} (\iconlg{layout-separator}) and \textbf{Stretches}
(\iconlg{layout-elastic-extender}) are invisible layout helpers.

% ── Chapter 23: Tab 3 --- Widgets ─────────────────────────────
\chapter{Paramètres --- Widgets}
\label{ch:settings-widgets}
\settingstab{3}{Widgets}

L'onglet Widgets gère la bibliothèque de définitions de widget de données utilisées par les deux barres.

\screenshot{settings-widgets}{Paramètres --- onglet Widgets avec liste de widgets et champs de définition}{}

\section{Champs de définition de widget}

\rowcolors{2}{LtGray}{white}
\begin{longtable}{L{2.8cm}L{10.9cm}}
\toprule
\rowcolor{Navy}
\color{white}\sffamily\bfseries Field &
\color{white}\sffamily\bfseries Purpose \\
\midrule
id            & Unique identifier used internally and in bar layout entries. \\
name          & Human-readable label shown in bars and this editor. \\
icon          & Optional SVG icon path. Accepts Qt resource paths or filesystem paths. \\
type          & Rendering format: \code{numerical}, \code{gauge}, \code{position},
                \code{datetime}, or \code{waypoint}. \\
signalKPath   & Full dot-separated Signal~K data path. \\
sourceUnit    & Unit in which Signal~K delivers the value (e.g.\ \code{m/s}, \code{rad}). \\
defaultUnit   & Unit to display in (e.g.\ \code{kn}, \code{deg}).
                Overridden by active unit preferences. \\
updatePolicy  & \code{instant} (default) or \code{fixed}. \\
period        & Target update period in milliseconds (default 1000). \\
minPeriod     & Minimum update period in milliseconds (default 200). \\
\bottomrule
\end{longtable}

% ── Chapter 24: Tab 4 --- Comfort ─────────────────────────────
\chapter{Paramètres --- Confort}
\label{ch:settings-comfort}
\settingstab{4}{Comfort}

L'onglet Comfort est l'éditeur de thème visuel pour les six préréglages de confort intégrés.

\screenshot{settings-comfort}{Paramètres --- Onglet confort : sélectionneur préréglé, montre couleur et section d'image de fond}{}

\section{Les Six Présets Confort}

\rowcolors{2}{LtGray}{white}
\begin{longtable}{C{1.5cm}L{2.0cm}L{4.2cm}L{6.0cm}}
\toprule
\rowcolor{Navy}
\color{white}\sffamily\bfseries Icon &
\color{white}\sffamily\bfseries Preset &
\color{white}\sffamily\bfseries Conditions &
\color{white}\sffamily\bfseries Key characteristics \\
\midrule
\iconlg{comfort-default} & Default & General fallback &
  Balanced mid-contrast. Usable in most conditions. \\
\iconlg{comfort-dawn}    & Dawn    & Around sunrise &
  Warm tones, slightly reduced brightness. Good contrast. \\
\iconlg{comfort-day}     & Day     & Full daylight, direct sun &
  Maximum brightness and contrast. For open-ocean sailing. \\
\iconlg{comfort-sunset}  & Sunset  & Around sunset &
  Warmer, transitional brightness. \\
\iconlg{comfort-dusk}    & Dusk    & Twilight &
  Reduced brightness, begins night-vision preservation. \\
\iconlg{comfort-night}   & Night   & Full dark, offshore watches &
Brillant minimum, fond sombre, accents rouges.
  Preserves night vision. Alarms remain high-contrast. \\
\bottomrule
\end{longtable}

\section{Boutons d'action}

\rowcolors{2}{LtGray}{white}
\begin{longtable}{C{1.5cm}L{3.0cm}L{9.2cm}}
\toprule
\rowcolor{Navy}
\color{white}\sffamily\bfseries Icon &
\color{white}\sffamily\bfseries Button &
\color{white}\sffamily\bfseries Action \\
\midrule
\iconlg{anchor-reset}     & Reset       &
  Clears all custom colour overrides, background images, and QSS for the selected preset. \\
\iconlg{delete-google}    & Reset All   &
  Clears customisation for all six presets. \\
\iconlg{route-export-iec} & Export QSS  &
  Saves the selected preset's effective stylesheet to a \code{.qss} file. \\
\iconlg{route-import-iec} & Import QSS  &
  Loads a \code{.qss} file as the custom stylesheet for the selected preset. \\
\bottomrule
\end{longtable}

\section{Propriétés de dépassement de couleur}

Dix propriétés de couleur sont personnalisables par préréglage:

\rowcolors{2}{LtGray}{white}
\begin{longtable}{L{3.6cm}L{10.1cm}}
\toprule
\rowcolor{Navy}
\color{white}\sffamily\bfseries Property &
\color{white}\sffamily\bfseries What it controls \\
\midrule
Window background  & Background colour of the Application Area and settings pages. \\
Text               & Primary text colour throughout the interface. \\
Button background  & Background of buttons and tool buttons in bars and drawers. \\
Button text        & Text and icon tint on buttons. \\
Border             & Borders, dividers, and frame outlines. \\
Accent             & Primary accent colour for highlights and selected states. \\
Accent text        & Text colour on accent-coloured backgrounds. \\
Icon               & Default tint for monochrome SVG icons. Set to red for Night preset to aid night vision. \\
Scroll track       & Background of scroll bar tracks. \\
Scroll knob        & Colour of the scroll bar handle. \\
\bottomrule
\end{longtable}

\begin{warningbox}
Ne jamais configurer un préréglage qui rend le texte d'alarme, les informations POB, ou la sécurité critique
lecture difficile à lire.
La lisibilité dans toutes les conditions prime sur l'esthétique.
\end{warningbox}

% ── Chapter 25: Tab 5 --- Connection ──────────────────────────
\chapter{Paramètres --- Connexion}
\label{ch:settings-connection}
\settingstab{5}{Connection}

\screenshot{settings-connection}{Paramètres --- Onglet de connexion : champ URL du serveur, ligne d'état, boutons d'action et journal de console}{}

\section{Champ URL du serveur}

An editable combo box pre-populated with the configured URL, \code{http://localhost:3000},
and \code{https://demo.signalk.org}.
FairWindSK normalises the URL: prepends \code{http://} if no scheme is present,
strips trailing slashes, and validates that only \code{http://} or \code{https://} are accepted.
Typed URLs are committed to the saved configuration only when you tap \key{Connect}.

\section{Ligne d'état}

Une seule ligne sous le champ URL affiche le jeton courant et l'état de connexion en utilisant
three dot-separated fields: \textbf{State}, \textbf{Permission}, and \textbf{Expires}.

\section{Boutons d'action}

\subsection{Connexion et pause}

Découpe la connexion Signal ~K.
\key{Connect} starts the connection; \key{Pause} suspends it without clearing the URL.

\subsection{Jeton de demande}

Initie le flux d'authentification de Signal~K read/write:

\begin{enumerate}
  \item FairWindSK posts a new access request to \skpath{/signalk/v1/access/requests}.
  \item On HTTP~202 + \code{PENDING}, FairWindSK opens the Signal~K admin interface in the
Tiroir de barre inférieur, navigué jusqu'à la sécurité ~> Page des demandes d'accès.
  \item FairWindSK polls the \code{href} every second until resolved.
  \item On \code{COMPLETED} with \code{APPROVED}, the token is stored in
        \code{fairwindsk.ini} and injected as the \code{JAUTHENTICATION} cookie in the
boutique de cookies du navigateur intégré.
\end{enumerate}

\subsection{Annuler}

Abandonne la demande de jeton en vol et dégage le href en attente.

\subsection{Supprimer le jeton}

Removes the stored token from memory, \code{fairwindsk.ini}, and the browser cookie store.
Déclenche un reconnect en mode public (non authentifié).

\section{Console Log}

Une zone de texte défilante en lecture seule sous les boutons affiche les événements de connexion :
messages de découverte mDNS, événements du cycle de vie de connexion, et résultats de la demande symbolique.

\section{mDNS Découverte automatique}

Sur les plates-formes de bureau, FairWindSK scanne le réseau local pour les serveurs Signal~K en utilisant
mDNS/Zeroconf sur quatre types de services:
\code{\_signalk-http.\_tcp}, \code{\_signalk-https.\_tcp}, \code{\_http.\_tcp},
and \code{\_https.\_tcp}.
Les serveurs découverts sont ajoutés automatiquement à la boîte combo URL.
mDNS est désactivé sur Android et iOS construit.

% ── Chapter 26: Tab 6 --- Signal K Paths ──────────────────────
\chapter{Paramètres --- Voies du signal K}
\label{ch:settings-signalk}
\settingstab{6}{Signal K Paths}

Carte les noms sémantiques que FairWindSK utilise en interne aux chemins de données actuels Signal~K
sur votre serveur.
Modifier n'importe quel champ pour changer où FairWindSK lit ou écrit cette valeur.
Les modifications prennent effet après un court délai de débonflement sans redémarrage nécessaire.

\screenshot{settings-signalk}{Paramètres --- Onglet Signal K Paths avec étiquette et champs de chemin modifiables}{}

\section{Cartographie des chemins clés}

\rowcolors{2}{LtGray}{white}
\begin{longtable}{L{3.0cm}L{5.6cm}L{5.1cm}}
\toprule
\rowcolor{Navy}
\color{white}\sffamily\bfseries Key &
\color{white}\sffamily\bfseries Default Signal~K path &
\color{white}\sffamily\bfseries Feature \\
\midrule
\code{pos}              & \skpath{navigation.position}                              & Top~Bar position, POB position \\
\code{sog}              & \skpath{navigation.speedOverGround}                       & Top~Bar SOG widget \\
\code{cog}              & \skpath{navigation.courseOverGroundTrue}                   & Top~Bar COG widget \\
\code{hdg}              & \skpath{navigation.headingTrue}                           & Top~Bar HDG widget \\
\code{dpt}              & \skpath{environment.depth.belowTransducer}                & Top~Bar depth widget \\
\code{notifications.pob}& \skpath{notifications.mob}                                & POB alarm \\
\code{anchor.position}  & \skpath{navigation.anchor.position}                      & Anchor Watch \\
\code{anchor.radius}    & \skpath{navigation.anchor.currentRadius}                 & Anchor Watch \\
\code{anchor.max}       & \skpath{navigation.anchor.maxRadius}                     & Anchor Watch \\
\code{anchor.rode}      & \skpath{winches.windlass.rode}                           & Anchor rode display \\
\code{anchor.actions.drop}   & \skpath{plugins.anchoralarm.dropAnchor}          & Anchor Drop button \\
\code{anchor.actions.raise}  & \skpath{plugins.anchoralarm.raiseAnchor}         & Anchor Raise button \\
\code{anchor.actions.radius} & \skpath{plugins.anchoralarm.setRadius}           & Anchor Set Radius button \\
\code{anchor.actions.up}     & \skpath{plugins.windlassctl.up}                  & Anchor Up button \\
\code{anchor.actions.down}   & \skpath{plugins.windlassctl.down}                & Anchor Down button \\
\code{environment.sun}  & \skpath{environment.sun}                                  & Auto comfort preset switching \\
\code{autopilot.state}  & \skpath{steering.autopilot.state}                        & Autopilot bar mode display \\
\code{autopilot.mode}   & \skpath{steering.autopilot.mode}                         & Autopilot engage commands \\
\code{rsa}              & \skpath{steering.rudderAngle}                            & Autopilot rudder indicator \\
\bottomrule
\end{longtable}

% ── Chapter 27: Tab 7 --- Units ────────────────────────────────
\chapter{Paramètres --- Unités}
\label{ch:settings-units}
\settingstab{7}{Units}

\screenshot{settings-units}{Paramètres --- Onglet Unités : nom prédéfini du serveur et lignes de remplacement par catégorie}{}

\section{Comment fonctionnent les préférences de l'unité}

FairWindSK lit les préférences de l'unité à partir du serveur Signal~K connecté.
Chaque ligne de catégorie montre : l'étiquette de catégorie, une boîte combo de remplacement (vos préférences locales),
et la valeur actuelle du serveur (lecture seule).
Si vous choisissez la même unité que le serveur, aucune redéfinition n'est stockée.
If different, it is saved in \code{unitPreferences.overrides.CATEGORY} in
\code{fairwindsk.json} and takes precedence.

% ── Chapter 28: Tab 8 --- Applications ────────────────────────
\chapter{Paramètres --- Demandes}
\label{ch:settings-apps}
\settingstab{8}{Applications}

Un éditeur à deux volets : le volet gauche gère les pages et dossiers du lanceur ; le volet droit est
la palette d'applications.

\screenshot{settings-apps}{Paramètres --- Onglet Applications : arborescence de page (gauche) et palette d'applications (droite)}{}

\section{Page Actions de l'arbre}

\rowcolors{2}{LtGray}{white}
\begin{longtable}{C{1.5cm}L{3.8cm}L{8.4cm}}
\toprule
\rowcolor{Navy}
\color{white}\sffamily\bfseries Icon &
\color{white}\sffamily\bfseries Action &
\color{white}\sffamily\bfseries What it does \\
\midrule
\iconlg{add-new-page}                    & Add page       & Creates a new empty launcher page. \\
\iconlg{show-page-details}               & Page details   & Opens the page details drawer (name, icon). \\
\iconlg{add-all-apps-to-current-page}    & Assign all apps & Adds all palette apps to the current page. \\
\iconlg{remove-all-apps-from-current-page}& Clear page    & Removes all tiles from the current page. \\
\iconlg{remove-all-apps-from-all-pages}  & Clear all pages & Removes all tiles from every page. \\
\bottomrule
\end{longtable}

\section{Palette Actions}

\rowcolors{2}{LtGray}{white}
\begin{longtable}{C{1.5cm}L{4.0cm}L{8.2cm}}
\toprule
\rowcolor{Navy}
\color{white}\sffamily\bfseries Icon &
\color{white}\sffamily\bfseries Action &
\color{white}\sffamily\bfseries What it does \\
\midrule
\iconlg{define-new-application-in-palette}  & Create app    & Opens the App Details drawer to create a new custom application. \\
\iconlg{show-details-selected-app-in-palette} & App details & Opens the App Details drawer to edit the selected application. \\
\iconlg{remove-selected-app-from-palette}   & Remove app    & Permanently removes the selected application from the palette. \\
\iconlg{add-selected-app-to-current-page}   & Add to page   & Adds the selected palette app to the current launcher page. \\
\bottomrule
\end{longtable}

\section{Détails de l'application Champs}

\rowcolors{2}{LtGray}{white}
\begin{longtable}{L{2.8cm}L{10.9cm}}
\toprule
\rowcolor{Navy}
\color{white}\sffamily\bfseries Field &
\color{white}\sffamily\bfseries Description \\
\midrule
Display name    & Name shown on the launcher tile and in the Top Bar when the app is active. \\
URL             & Full URL (\code{http://} or \code{https://}). \\
Icon path       & Qt resource path, filesystem path, or bundled icon path. \\
Web zoom level  & Percentage (default 100\%). Increase for larger buttons; decrease for more content. \\
Active          & Whether the tile is visible on the launcher. \\
Order           & Secondary sort hint within a page. \\
\bottomrule
\end{longtable}

% ── Chapter 29: Tab 9 --- System ──────────────────────────────
\chapter{Paramètres --- Système}
\label{ch:settings-system}
\settingstab{9}{System}

The System tab is divided into two panes: \textbf{Performance} (left) and
\textbf{Diagnostics and Logging} (right).

\screenshot{settings-system-perf}{Paramètres --- Onglet système Panneau de performance : barres mémoire, processeur par cœur, réseau, métriques Raspberry Pi}{}

\section{Panneau de performance}

\subsection{Lecture de la mémoire}

\begin{itemize}
  \item \textbf{Process memory} --- resident RAM used by FairWindSK
        (from \code{/proc/self/statm} on Linux; Mach task API on macOS).
  \item \textbf{Total memory} --- installed RAM
        (from \code{/proc/meminfo} on Linux; \code{sysctl hw.memsize} on macOS).
  \item \textbf{Memory usage bar} --- process memory as a percentage of total RAM.
\end{itemize}

\subsection{Chargement du processeur}

Une rangée de barres de progression horizontales, une par noyau du processeur, mise à jour chaque seconde.
On Linux, data is read from \code{/proc/stat};
on macOS from the Mach \code{host\_processor\_info} API.

\subsection{Adresses réseau}

Liste les adresses IPv4 de toutes les interfaces réseau actives et non loopback, mises à jour toutes les 10 secondes.
Utile pour confirmer quelle adresse IP d'autres appareils devraient utiliser pour atteindre cette installation.

\subsection{Diagnostics de la framboise Pi}

Quand le serveur Signal~K connecté expose les métriques Raspberry~Pi
(via the \code{signalk-server-raspberrypi} plugin on path \skpath{environment.rpi}),
une boîte de groupe Raspberry~Pi apparaît montrant :

\rowcolors{2}{LtGray}{white}
\begin{longtable}{L{3.6cm}L{10.1cm}}
\toprule
\rowcolor{Navy}
\color{white}\sffamily\bfseries Metric &
\color{white}\sffamily\bfseries Description \\
\midrule
CPU temperature   & SoC CPU temperature in \textdegree{}C. Above 80\textdegree{}C indicates thermal throttling. \\
GPU temperature   & VideoCore GPU temperature in \textdegree{}C. \\
CPU utilisation   & Overall CPU load as a percentage. \\
Memory utilisation& RAM usage as a percentage. \\
SD utilisation    & Storage usage on the SD card as a percentage. \\
\bottomrule
\end{longtable}

\section{Diagnostics et couverture d'exploitation}

\screenshot{settings-system-log}{Paramètres --- Onglet système Diagnostics volet: niveau de log, logs persistants, historique d'interaction, chemins, Affichage des journaux et Affichage des boutons Rapports}{}

\subsection{Niveau du journal}

\rowcolors{2}{LtGray}{white}
\begin{longtable}{L{2.4cm}L{4.6cm}L{6.7cm}}
\toprule
\rowcolor{Navy}
\color{white}\sffamily\bfseries Level &
\color{white}\sffamily\bfseries Messages included &
\color{white}\sffamily\bfseries When to use \\
\midrule
Off (default) & No messages. & Normal operation. Minimises SD card writes on Raspberry~Pi. \\
Critical      & Fatal errors only. & Minimal logging with crash detection. \\
Warning       & Critical + non-fatal anomalies. & Initial installation testing. \\
Info          & Warning + lifecycle events. & Commissioning and setup. \\
Debug         & Info + detailed internal state. & Troubleshooting. Significant overhead. \\
Full          & All of the above + verbose tracing. & Developer use only. Very high overhead. \\
\bottomrule
\end{longtable}

\subsection{Registres persistants}

Une fois activé, FairWindSK écrit un journal complet des messages pour chaque exécution à l'utilisateur
répertoire de données d'application, avec le processus de démarrage timestamp dans le nom de fichier.

\subsection{Historique des interactions}

Lorsque activé, FairWindSK écrit un journal léger de navigation et de contrôle de l'opérateur
actions à côté de chaque journal d'exécution.
Utile pour l'examen post-incident.

\subsection{Afficher les journaux et les rapports Boutons}

Ouvrez l'explorateur de journaux dans le tiroir horizontal de la barre inférieure.
Les fichiers journaux apparaissent à gauche; les entrées de journaux parsées apparaissent à droite.
Les paquets de rapports de crash (créés automatiquement lorsque FairWindSK n'est pas sorti proprement) sont
accessible via \key{View Reports}.

\section{Boutons de gestion de configuration}

\rowcolors{2}{LtGray}{white}
\begin{longtable}{L{3.2cm}L{10.5cm}}
\toprule
\rowcolor{Navy}
\color{white}\sffamily\bfseries Button &
\color{white}\sffamily\bfseries Action \\
\midrule
Réinitialiser&
Élimine les modifications et recharge la configuration active.
  FairWindSK keeps running. \\
Restaurer les défauts&
Remplace tous les paramètres par les paquets par défaut.
  All customisation is lost. Confirmation required. \\
Importation&
  Imports a \code{fairwindsk.json} file from the filesystem.
  Validates, atomically replaces the active configuration, and reloads immediately. \\
Exportations&
  Copies the current \code{fairwindsk.json} to a chosen location. \\
\bottomrule
\end{longtable}

\begin{warningbox}
\key{Restore Defaults} is irreversible without a backup.
Exportez d'abord votre configuration si vous souhaitez la restaurer.
\end{warningbox}

\section{Boutons de contrôle du système}

\rowcolors{2}{LtGray}{white}
\begin{longtable}{L{3.2cm}L{10.5cm}}
\toprule
\rowcolor{Navy}
\color{white}\sffamily\bfseries Button &
\color{white}\sffamily\bfseries Action \\
\midrule
Redémarrer FairWindSK&
Enregistre les paramètres, puis ferme FairWindSK avec le code de sortie ~1.
  The Raspberry~Pi~OS startup helper relaunches on exit code~1. \\
Quitter FairWindSK&
Enregistre les paramètres et ferme proprement (sortie code ~0).
  The startup helper does \textbf{not} relaunch after a clean exit. \\
\bottomrule
\end{longtable}

\begin{warningbox}
Quitter FairWindSK sur un écran de barre dédié laisse l'affichage vide.
Assurez-vous qu'un autre membre de l'équipage garde une surveillance sécuritaire avant de quitter le service.
\end{warningbox}
